\documentclass{article}

\usepackage[english]{babel}
\usepackage[
    letterpaper,
    top=2cm,
    bottom=2cm,
    left=3cm,
    right=3cm,
    marginparwidth=1.75cm
]{geometry}
\usepackage{amsmath}
\usepackage{amssymb}
\usepackage[ruled,vlined,linesnumbered]{algorithm2e}
\SetKw{Return}{return}
\usepackage{booktabs}
\usepackage[colorlinks=true,allcolors=blue]{hyperref}

\title{Workflow Lifecycle Scheduling Fragment}
\author{}
\date{}

\begin{document}
\maketitle

\section{Core Idea and Candidate Contributions}

The central design in this fragment is to make an ordered model-lifecycle plan,
rather than a ready activation alone, the unit of scheduling. A readiness-only
workflow scheduler chooses placement after an activation becomes ready. The
proposed scheduler uses predicted DAG readiness and resource demand to act
earlier: it reuses a compatible resident model, prefetches a missing model
before its activation is released, evicts low-value idle models when capacity
is scarce, or defers a transition whose complete path is unsafe.

The resulting mechanism makes three coupled decisions:

\begin{itemize}
    \item \textbf{Lifecycle plans as scheduling units.} Each decision selects an
    activation together with an ordered sequence of reuse, load, eviction,
    prefetch, and dispatch actions.
    \item \textbf{Prediction-guided reuse and prefetch.} Cross-session DAG
    lookahead retains models with imminent reuse and overlaps model loading with
    predecessor execution.
    \item \textbf{Value-aware eviction with path safety.} Eviction accounts for
    the reload and lost-reuse cost of each victim, while admission checks the
    peak memory of the entire transition rather than its final state alone.
\end{itemize}

\section{Motivation}

\subsection{Chain Workflow}

A chain workflow with $n$ sequential LLM nodes and $m$ independent sessions
exposes a structural resource-waste problem. Let $t_{i,j}$ be the active LLM
execution time of node $i$ on session $j$, and let
$A=\sum_i\sum_j t_{i,j}$ be the active GPU-seconds. Let $G\ge 0$ denote the
remaining execution span in which none of the chain models is active, including
arrival gaps and intervening non-LLM work. The makespan is $T=A+G$. If all $n$
model instances remain resident throughout this span, their resident GPU-seconds
are $R_{\mathrm{static}}=nT$, yielding
\[
    \rho_{\mathrm{static}}
    = \frac{R_{\mathrm{static}}}{A}
    = n\left(1+\frac{G}{A}\right)
    \ge n.
\]

We validate this bound on MBPP~\cite{austin2021program} with a
\texttt{coder $\rightarrow$ reviewer $\rightarrow$ repair} chain. The three
stages use Qwen3-14B, Qwen3-8B, and Qwen3-14B on separate GPUs. Across 60
sessions, active and resident GPU-seconds are 5,656 and 17,017, respectively,
giving a measured gap of 3.009 (Table~\ref{tab:mbpp_empirical}).

\begin{table}[htbp]
    \centering
    \caption{MBPP chain resource gap ($n=3$, $m=60$).}
    \label{tab:mbpp_empirical}
    \begin{tabular}{lr}
        \toprule
        Metric & Value \\
        \midrule
        Active GPU-seconds & 5,656 \\
        Resident GPU-seconds & 17,017 \\
        Resource gap $\rho$ & 3.009 \\
        \bottomrule
    \end{tabular}
\end{table}

This bound is specific to a chain whose distinct stage models are resident
concurrently. It should not be generalized to workflows backed by one shared
model or to runtimes that already load models on demand.

\section{Lifecycle Scheduling Design}
\label{sec:lifecycle-scheduling-design}

At each scheduling event, the scheduler chooses a pair $(x,b)$: an activation
$x$ and the lifecycle plan $b$ that makes its model state available. This joint
decision is the core design choice. Reuse changes the cost of dispatch,
prefetch changes when memory is committed, and eviction changes which future
activations will pay a reload; treating them as independent policies loses
these interactions.

\subsection{Scheduling Model}

The runtime serves multiple workflow instances from a shared accelerator and
model pool. An activation $x=(w,s,v)$ identifies workflow $w$, session $s$, and
node $v$; $M_x$ is its required model. Workflow-local state remains isolated,
while a compatible physical model instance may be reused across activations and
sessions. For each device $g$, $\mathcal{Z}_t(g)$ partitions resident instances
into active, reserved, and reclaimable states. Only reclaimable instances may be
evicted.

The scheduler operates on two frontiers: ready activations
$\mathcal{R}_t$, whose dependencies are satisfied, and near-ready activations
$\mathcal{N}_t$, whose release times can be estimated from running
predecessors. Together, these frontiers reveal both immediate demand and the
next compatible use of each resident model.

\subsection{Prediction-Guided Lifecycle Plans}

For activation $x$ on device $g$, the forecaster exposes
\begin{equation}
    \widehat{K}(x,g)=\left\langle
        \hat{t}_{\mathrm{ready}},
        \hat{t}_{\mathrm{load}},
        \hat{t}_{\mathrm{run}},
        U_{\mathrm{mem}}
    \right\rangle,
    \label{eq:lifecycle-contract}
\end{equation}
where the time estimates guide ordering and prefetch, while
$U_{\mathrm{mem}}$ is a calibrated upper bound used for admission. The GNN may
provide execution time and peak-memory estimates from the model graph and token
regime; online measurements provide model-load costs and current residency.
The frontiers additionally induce a reuse value $\widehat{V}_t(m)$ for resident
instance $m$, representing the load latency and workflow delay avoided by
retaining it for predicted future activations.

For each activation-device pair, the planner constructs semantically distinct
action sequences:
\begin{equation}
    \begin{aligned}
        \Pi_{\mathrm{ready}}(x,g)=\{&
        \langle\mathsf{Reuse}(m)\rangle,
        \langle\mathsf{Load}(M_x,g)\rangle,\\
        &\langle\mathsf{Evict}(E),\mathsf{Load}(M_x,g)\rangle\},\\
        \Pi_{\mathrm{near}}(x,g)=\{&
        \langle\mathsf{Prefetch}(M_x,g)\rangle,\\
        &\langle\mathsf{Evict}(E),\mathsf{Prefetch}(M_x,g)\rangle\}.
    \end{aligned}
    \label{eq:lifecycle-plans}
\end{equation}
Here, $m$ must be a compatible resident instance and
$E\subseteq\mathcal{Z}^{\mathrm{reclaim}}_t(g)$ is a victim set. Reuse avoids a
cold load: an idle compatible instance yields a \textsc{Reuse} plan with no
model-transition latency or additional weight residency. Proactive model
preloading, denoted as \textsc{Prefetch}, begins before dependencies release the
activation. Each successful preparation plan ends in \textsc{Dispatch} after
dependencies are satisfied. Eviction is coupled to the load or prefetch that
requires its reclaimed capacity. \textsc{Defer} remains a first-class plan when
no transition is safe.

A near-ready activation becomes eligible for prefetch when
\begin{equation}
    \hat{t}_{\mathrm{ready}}(x)-t
    \le \hat{t}_{\mathrm{load}}(x)+\delta,
    \label{eq:prefetch-window}
\end{equation}
where $\delta$ absorbs forecast uncertainty and $\hat{t}_{\mathrm{load}}(x)$ is
the shortest predicted load time among candidate-device plans. Ready work always
precedes speculative prefetch.

Eviction is prediction-guided rather than a fixed recency rule. The scheduler
selects the lowest-regret victim set that can supply the required capacity:
\begin{equation}
    \begin{aligned}
        \widehat{C}_{\mathrm{evict}}(E)
        &=\sum_{m\in E}
        \left(\widehat{C}_{\mathrm{reload}}(m)+\widehat{V}_t(m)\right)\\
        E^*(x,g)&=\arg\min_E\widehat{C}_{\mathrm{evict}}(E)\\
        \text{s.t.}\quad
        &E\subseteq\mathcal{Z}^{\mathrm{reclaim}}_t(g),\\
        &M_{\mathrm{free}}(g)+\sum_{m\in E}M(m)
        \ge U_{\mathrm{mem}}(x,g).
    \end{aligned}
    \label{eq:eviction-regret}
\end{equation}
Here, $M(m)$ is the reclaimable memory footprint of instance $m$. Thus,
eviction trades reclaimed capacity against the expected cost of reloading models
and destroying near-future reuse.

\subsection{Plan Scoring and Admission}

The scheduler ranks feasible plans using
\begin{equation}
    \widehat{J}(b)
    =\widehat{\Delta\mathrm{JCT}}(b)
    +\lambda_e\widehat{C}_{\mathrm{evict}}(E_b)
    +\lambda_u\widehat{C}_{\mathrm{uncertainty}}(b),
    \label{eq:lifecycle-objective}
\end{equation}
where $E_b=\emptyset$ for plans without eviction and $\lambda_e,\lambda_u$
control eviction-regret and uncertainty penalties.
$\widehat{\Delta\mathrm{JCT}}$ includes execution and queueing delay, but counts
only the portion of model loading exposed on the predicted workflow critical
path. Direct reuse therefore pays no load cost, successful prefetch hides it,
and eviction still pays its expected future regret.

A plan is admitted only if dependencies and model compatibility hold and its
calibrated path bound satisfies
\begin{equation}
    U_{\mathrm{path}}(b,g,\tau)\le M_g
    \quad \forall \tau\in b,
    \label{eq:path-memory-safety}
\end{equation}
where $M_g$ is device capacity and the bound includes protected resident state,
the incoming model, execution workspace, and transient overlap before evicted
state is fully reclaimed.
Selecting a plan reserves every model state and capacity interval on which this
check depends.

\subsection{Scheduling Procedure}

The runtime invokes Algorithm~\ref{alg:lifecycle-scheduling-event} on arrivals,
completions, and material state changes. Each invocation admits ready work first,
considers prefetch only when ready work cannot progress, and defers when no safe
plan exists.

\begin{algorithm}[t]
    \caption{Prediction-guided model-lifecycle scheduling.}
    \label{alg:lifecycle-scheduling-event}
    \KwIn{Ready frontier $\mathcal{R}_t$, near-ready frontier $\mathcal{N}_t$,
    runtime state $\mathcal{S}_t$, forecaster $f_\theta$}
    \KwOut{A reserved lifecycle plan or \textsc{Defer}}

    $\mathcal{R} \gets \operatorname{RankReady}(\mathcal{R}_t,\mathcal{S}_t)$\;
    \ForEach{$x\in\mathcal{R}$}{
        $\mathcal{P}\gets\operatorname{ReusePlans}(x,\mathcal{S}_t)$\;
        $\mathcal{P}\gets\mathcal{P}\cup
        \operatorname{LoadPlans}(x,\mathcal{S}_t)$\;
        $\mathcal{P}\gets\mathcal{P}\cup
        \operatorname{EvictLoadPlans}(x,\mathcal{S}_t)$\;
        $\mathcal{B}\gets\operatorname{SafePaths}(\mathcal{P},f_\theta)$\;
        \If{$\mathcal{B}\neq\emptyset$}{
            $b^*\gets\arg\min_{b\in\mathcal{B}}\widehat{J}(b)$\;
            \Return{$\operatorname{Reserve}(b^*,\mathcal{S}_t)$}\;
        }
    }
    $\mathcal{N}\gets
    \operatorname{PrefetchEligible}(\mathcal{N}_t,f_\theta,t,\delta)$\;
    \ForEach{$x\in\operatorname{RankNear}(\mathcal{N},\mathcal{S}_t)$}{
        $\mathcal{P}\gets\operatorname{PrefetchPlans}(x,\mathcal{S}_t)$\;
        $\mathcal{P}\gets\mathcal{P}\cup
        \operatorname{EvictPrefetchPlans}(x,\mathcal{S}_t)$\;
        $\mathcal{B}\gets\operatorname{SafePaths}(\mathcal{P},f_\theta)$\;
        \If{$\mathcal{B}\neq\emptyset$}{
            $b^*\gets\arg\min_{b\in\mathcal{B}}\widehat{J}(b)$\;
            \Return{$\operatorname{Reserve}(b^*,\mathcal{S}_t)$}\;
        }
    }
    \Return{\textsc{Defer}}\;
\end{algorithm}

\textsc{RankReady} remains the interface to the upstream workflow policy.
The lifecycle layer owns model reuse, prefetch timing, victim valuation, and
path admission. The plan constructors expose only the semantic transitions in
Eq.~\ref{eq:lifecycle-plans}; engine-specific bookkeeping remains outside the
algorithm.

\section{Related Work}

\paragraph{Workflow-aware LLM serving.} Parrot exposes request dependencies and
prompt-prefix relationships to support
application-centric placement and reuse~\cite{lin2024parrot}. Kairos ranks
requests using workflow-derived remaining-latency distributions and dispatches
them using memory demand~\cite{chen2025kairos}. Murakkab separates declarative
workflow structure from execution configuration and uses profiles to choose
models and hardware under SLOs~\cite{chaudhry2025murakkab}. SAGA treats an agent
workflow as the scheduling unit to preserve KV continuity and provide
workflow-level fairness~\cite{guo2026saga}.

Together, these systems demonstrate the value of workflow visibility at the
orchestration and serving layers. The contribution of this fragment is to make
the predicted model-state transition itself schedulable: every activation is
admitted together with a ranked reuse, prefetch, or eviction plan whose complete
transition is resource-safe.

\bibliographystyle{alpha}
\bibliography{sample}

\end{document}
